\documentclass[12pt,a4paper]{article}
\usepackage[utf8]{inputenc}
\usepackage{amsmath,amssymb,amsthm}
\usepackage{algorithm}
\usepackage{algorithmic}
\usepackage{geometry}
\usepackage{booktabs}
\geometry{margin=1in}

\newtheorem{theorem}{Theorem}
\newtheorem{lemma}[theorem]{Lemma}
\newtheorem{proposition}[theorem]{Proposition}
\theoremstyle{definition}
\newtheorem{definition}{Definition}
\theoremstyle{remark}
\newtheorem*{remark}{Remark}

\title{Project Euler Problem 24: Lexicographic Permutations}
\author{}
\date{}

\begin{document}
\maketitle

\section{Problem Statement}

Determine the millionth lexicographic permutation of the digits $\{0, 1, 2, 3, 4, 5, 6, 7, 8, 9\}$.

\section{Definitions}

\begin{definition}[Factorial Number System]
The \emph{factoradic} represents a non-negative integer $k$ as
\[
k = \sum_{i=1}^{m} d_i \cdot i!, \qquad 0 \leq d_i \leq i.
\]
\end{definition}

\begin{definition}[Lehmer Code]
Given a permutation $\pi$ of $\{0, \ldots, n-1\}$, its \emph{Lehmer code} is $(l_0, \ldots, l_{n-1})$ where $l_i = |\{j > i : \pi(j) < \pi(i)\}|$.
\end{definition}

\section{Theorems and Proofs}

\begin{theorem}[Factoradic Uniqueness]
Every integer $k$ with $0 \leq k < n!$ has a unique factoradic representation $(d_{n-1}, \ldots, d_1)$ with $0 \leq d_i \leq i$.
\end{theorem}

\begin{proof}
By induction on $n$. For $n = 1$, $k = 0$ trivially. For the inductive step, write $k = d_{n-1} \cdot (n-1)! + r$ via the division algorithm. Since $k < n!$, we have $0 \leq d_{n-1} \leq n-1$. Since $0 \leq r < (n-1)!$, the inductive hypothesis gives a unique representation of $r$. Uniqueness of the full representation follows from uniqueness of quotient and remainder.
\end{proof}

\begin{theorem}[Lehmer Code Bijection]
The $k$-th permutation (0-indexed) of $n$ elements in lexicographic order is constructed from the factoradic of~$k$: at step~$i$, select the element of rank $d_{n-1-i}$ among remaining elements.
\end{theorem}

\begin{proof}
By induction. The $n!$ permutations partition into $n$ blocks of $(n-1)!$, where block~$j$ starts with the $j$-th smallest element. The first digit is $d_{n-1} = \lfloor k/(n-1)!\rfloor$, and the remainder $r = k \bmod (n-1)!$ indexes into the sub-permutation of the remaining $n-1$ elements. The inductive hypothesis completes the construction.
\end{proof}

\begin{lemma}[Feasibility]
$k = 999999 < 10! = 3628800$, so the factoradic representation exists with $n = 10$.
\end{lemma}

\section{Computation}

\begin{center}
\begin{tabular}{@{}cccccc@{}}
\toprule
Step $i$ & $k$ & $(9{-}i)!$ & $d_{9-i}$ & Available & Selected \\
\midrule
0 & 999999 & 362880 & 2 & 0123456789 & 2 \\
1 & 274239 & 40320 & 6 & 013456789 & 7 \\
2 & 32319 & 5040 & 6 & 01345689 & 8 \\
3 & 2079 & 720 & 2 & 0134569 & 3 \\
4 & 639 & 120 & 5 & 014569 & 9 \\
5 & 39 & 24 & 1 & 01456 & 1 \\
6 & 15 & 6 & 2 & 0456 & 5 \\
7 & 3 & 2 & 1 & 046 & 4 \\
8 & 1 & 1 & 1 & 06 & 6 \\
9 & 0 & 1 & 0 & 0 & 0 \\
\bottomrule
\end{tabular}
\end{center}

\section{Algorithm}

\begin{algorithm}[H]
\caption{$k$-th Lexicographic Permutation}
\begin{algorithmic}[1]
\REQUIRE Sorted set $D$ of $n$ elements, 1-based index $k$
\ENSURE The $k$-th lexicographic permutation of $D$
\STATE $k \gets k - 1$ \COMMENT{convert to 0-indexed}
\STATE $\mathrm{result} \gets [\ ]$
\FOR{$i = 0$ \TO $n - 1$}
    \STATE $f \gets (n - 1 - i)!$
    \STATE $q \gets \lfloor k / f \rfloor$; \quad $k \gets k \bmod f$
    \STATE Append $D[q]$ to $\mathrm{result}$; remove $D[q]$ from $D$
\ENDFOR
\RETURN $\mathrm{result}$
\end{algorithmic}
\end{algorithm}

\section{Complexity Analysis}

\begin{proposition}
The algorithm runs in $O(n^2)$ time and $O(n)$ space.
\end{proposition}

\begin{proof}
The loop iterates $n$ times. Each iteration removes an element from a list of size $O(n)$, costing $O(n)$. All other per-iteration work is $O(1)$. Total: $O(n^2)$. Space: two lists of size $O(n)$.
\end{proof}

\begin{remark}
Using an order-statistic tree, select-and-delete costs $O(\log n)$ per step, reducing total time to $O(n \log n)$.
\end{remark}

\section{Answer}

\[
\boxed{2783915460}
\]

\end{document}
